\chapter{Conclusão}\label{chp:CONCLUSAO}

Este capítulo encerra o trabalho com um breve resumo do que foi apresentado, considerações sobre o que foi alcançado em relação ao objetivo proposto na Introdução (Capítulo~\ref{chp:INTRODUCAO}), as limitações identificadas e sugestões para trabalhos futuros.

\section{Considerações finais}

Este trabalho partiu da necessidade de facilitar a busca por atas de reuniões da UFRRJ, propondo um chatbot de regras capaz de mapear as perguntas de discentes, docentes e do público acadêmico em geral para as atas mais adequadas ao conteúdo de suas buscas. Para isso, o Capítulo~\ref{chp:FUNDAMENTACAO} situou o problema dentro da literatura sobre chatbots e sobre Busca e Recuperação da Informação (BRI), apresentando a taxonomia dos principais tipos de chatbot disponíveis, o comparativo entre essas opções e os requisitos do projeto, e os modelos clássicos de recuperação, com destaque para o modelo probabilístico e o algoritmo BM25, utilizado pelo Elasticsearch para o ranqueamento padrão de resultados.

A partir dessa fundamentação, o Capítulo~\ref{chp:PROPOSTA} caracterizou o problema central deste trabalho: a desconexão entre o vocabulário livre empregado pelo usuário em uma conversa e a forma como o conteúdo das atas está organizado no índice do Elasticsearch. A solução proposta foi um roteiro guiado de diálogo, que conduz o usuário por um conjunto fechado de critérios de busca, traduz suas respostas em parâmetros estruturados e delimita, de forma explícita, o que pode ou não ser buscado, fechando essa lacuna sem exigir que o usuário conheça a organização do índice ou a sintaxe de consulta do motor de busca.

O Capítulo~\ref{chp:EXPERIMENTOS} descreveu a implementação dessa proposta, composta pelo Botpress, responsável pelo fluxo de diálogo; pelo indexador, desenvolvido em Node.js com Express, responsável por extrair o conteúdo das atas em PDF e montar a consulta ao Elasticsearch; pelo Elasticsearch propriamente dito, responsável pelo armazenamento e pelo ranqueamento via BM25; e por um bot do Discord, oferecido como canal alternativo de interação. Foi detalhado o mapeamento de cada um dos sete critérios de busca do roteiro (departamento, disciplina, docentes, discentes, período, processo/edital e progressão), além do campo de outros assuntos, para cláusulas da consulta em \textit{Query DSL}, com pesos distintos conforme a especificidade de cada critério, e a aplicação do score para reordenar os três melhores resultados antes de sua exibição ao usuário.

Em relação ao objetivo definido na Introdução, definido por disponibilizar uma ferramenta conversacional capaz de interpretar perguntas em linguagem natural e mapeá-las para as atas institucionais mais pertinentes, o trabalho alcançou esse resultado por meio da combinação de um chatbot baseado em regras com um roteiro guiado de diálogo, com as definições dos paramêtros que seriam oferecidos para o usuário realizar a personalização da busca. O comparativo do Capítulo~\ref{chp:FUNDAMENTACAO} (Tabela~\ref{tab:comparativo-chatbots}) evidenciou que a escolha por um chatbot baseado em regras é coerente com o escopo fechado do problema, com a ausência de orçamento para licenciamento de ferramentas e com a necessidade de integração via API com o Elasticsearch. O roteiro guiado de diálogo, por sua vez, mostrou-se uma solução viável para o problema de geração de consultas a partir de uma interação conversacional, ao definir de forma explícita quais critérios podem ser buscados e como cada resposta do usuário se traduz em um componente da consulta, sem exigir do usuário qualquer conhecimento sobre a organização do índice ou sobre o BM25.

\section{Limitações e trabalhos futuros}

Por se tratar de um chatbot baseado em regras, o sistema opera sobre um vocabulário e um conjunto de critérios de busca fechados e definidos previamente no roteiro de diálogo. Conforme discutido nos Capítulos~\ref{chp:FUNDAMENTACAO} e \ref{chp:PROPOSTA}, essa característica é adequada ao escopo do problema, mas limita a interpretação de perguntas que não se encaixem nos sete critérios previstos (departamento, disciplina, docentes, discentes, período, processo/edital e progressão) ou no campo de assunto geral, o qual concentra qualquer informação fornecida pelo usuário que não corresponda a esses critérios ou a um período de datas reconhecível. Como trabalho futuro, sugere-se o mapeamento de novos critérios de busca além dos sete atualmente previstos, a partir do levantamento de outros aspectos recorrentes no conteúdo das atas institucionais que ainda não tenham sido contemplados pelo roteiro --- por exemplo, tipos específicos de deliberação ou categorias de assunto frequentemente tratadas pelo colegiado --- de modo a reduzir, ao longo do tempo, a dependência do campo de assunto geral para temas que poderiam ser tratados como critérios próprios, com peso e tipo de correspondência específicos na consulta ao Elasticsearch.


Uma extensão mais ampla seria a evolução do chatbot atual, baseado exclusivamente em regras, para um modelo híbrido, que combine o roteiro guiado já implementado com uma camada de IA generativa para a interpretação de perguntas em linguagem mais livre, sobretudo naquelas que hoje são apenas direcionadas ao campo de assunto geral. Uma implementação desse tipo, no entanto, precisaria de um cuidado adicional para mitigar o risco de alucinação já discutido na Seção~\ref{sec:problema-gerar-consulta}: a geração de respostas pelo componente de IA deveria ser estritamente fundamentada no conteúdo efetivamente recuperado do Elasticsearch (uma arquitetura de \textit{Retrieval-Augmented Generation}, ou RAG), com instruções de contexto bem definidas que restrinjam o modelo a responder apenas com base nas atas retornadas pela busca, evitando que o chatbot produza informações não embasadas nos documentos institucionais, como discutido no exemplo do chatbot da concessionária Chevrolet de Watsonville \cite{aiincidentdb622, venturebeat-chevy}.

O sistema também está, atualmente, restrito ao Botpress e ao bot do Discord como canais de interação. Como trabalho futuro, é possível integrar o mesmo roteiro de diálogo a outras plataformas de uso comum pela comunidade acadêmica, como um \textit{widget} de chat no próprio site da universidade ou do departamento, ou um bot para o Telegram, ampliando o alcance do sistema sem a necessidade de duplicar a lógica de busca já implementada no indexador, dado que o Botpress permite a integração com múltiplos canais de comunicação a partir de um único fluxo de diálogo.

Por fim, no que diz respeito ao armazenamento dos documentos, as atas em PDF utilizadas neste trabalho são mantidas localmente, em uma pasta no mesmo ambiente onde o indexador é executado, sendo apenas seu conteúdo textual extraído e enviado ao Elasticsearch. Essa abordagem é suficiente para o escopo deste trabalho, mas limita a escalabilidade e a disponibilidade do sistema a um único ambiente local. Como trabalho futuro, sugere-se a migração do armazenamento das atas para um serviço de armazenamento em nuvem ou para um banco de dados próprio para arquivos, permitindo que o indexador e a rota de acesso aos arquivos PDF (\texttt{/arquivos}) passem a consultar as atas a partir desse serviço externo, em vez de um diretório local, o que facilitaria tanto a hospedagem do sistema em um ambiente de produção quanto a atualização do acervo de atas por outros membros do departamento.

Sugere-se ainda, como trabalho futuro, a realização de uma avaliação de usabilidade com usuários reais do Departamento de Ciência da Computação da UFRRJ, envolvendo discentes, docentes e o público acadêmico em geral, de modo a verificar, em uso real, a efetividade do roteiro guiado na recuperação das atas desejadas e a identificar critérios de busca ou formulações de pergunta não previstos no escopo atual.
